% 矮行星（综述）
% license CCBYSA3
% type Wiki

本文根据 CC-BY-SA 协议转载翻译自维基百科\href{https://en.wikipedia.org/wiki/Dwarf_planet}{相关文章}。

\begin{figure}[ht]
\centering
\includegraphics[width=6cm]{./figures/d21611a5210b90a2.png}
\caption{} \label{fig_Dwarf_1}
\end{figure}
“矮行星是一类直接绕太阳运行、具有行星质量但体积较小的天体。它们的质量足以使自身达到由重力塑形的近球形，但不足以像太阳系中八大经典行星那样在轨道上实现支配地位。矮行星的原型是冥王星，它曾被视为一颗行星长达数十年，直到 2006 年‘矮行星’概念被正式采用。许多行星地质学家认为矮行星与具有行星质量的卫星都应被视为行星\(^\text{[1]}\)，但自 2006 年以来，国际天文学联合会（IAU）和许多天文学家已将它们排除在行星名录之外。

矮行星可能具有地质活动能力——这一预期在 2015 年“黎明号”对谷神星的探测和“新视野号”对冥王星的探测中得到证实。因此，它们成为行星地质学家特别关注的对象。

天文学家普遍同意至少有九个最大的候选体属于矮行星——按直径从大到小大致排列为：冥王星、厄里斯、妊神星、鸟神星、公公星、夸欧尔、塞德娜、谷神星和奥库斯。对于第十大候选体萨拉西娅仍存在相当的不确定性，因此可将其视为临界案例。在这十个天体中，已有两个被航天器访问过（冥王星和谷神星），另有七个至少拥有一颗已知卫星（厄里斯、妊神星、鸟神星、公公星、夸欧尔、奥库斯和萨拉西娅），这使得它们的质量得以确定，从而可估算其密度。质量和密度又可用于构建地球物理模型，以尝试判断这些世界的性质。只有塞德娜既未被探测器访问，也没有已知卫星，因此其质量难以准确估计。一些天文学家还将许多更小的天体也纳入矮行星的范畴\(^\text{[2]}\)，但目前并无共识认为这些小天体很可能是矮行星。”
\subsection{概念的历史}
\begin{figure}[ht]
\centering
\includegraphics[width=6cm]{./figures/0ce0512cdfe2e0c4.png}
\caption{冥王星及其卫星卡戎的近真彩色影像。二者的间距按比例绘制} \label{fig_Dwarf_2}
\end{figure}
\begin{figure}[ht]
\centering
\includegraphics[width=6cm]{./figures/6ae4d90716be2481.png}
\caption{4 号灶神星（Vesta），一度曾是矮行星的天体\(^\text{[3]}\)} \label{fig_Dwarf_3}
\end{figure}
自1801年起，天文学家发现了谷神星以及位于火星与木星之间的其他天体，这些天体在随后的数十年间都被视为行星。自那时直到大约1851年，当行星数量已达到23颗时，天文学家开始使用“asteroid”（来自希腊语，意为“类星的”或“星形的”）一词来指代较小的天体，并开始将它们区分为小行星而非大行星。\(^\text{[4]}\)

随着1930年冥王星的发现，大多数天文学家认为太阳系包含九颗主要行星，以及成千上万的显著更小的天体（小行星与彗星）。在将近50年的时间里，人们曾认为冥王星比水星更大，\(^\text{[5][6]}\)但在1978年冥王星的卫星卡戎被发现后，天文学家得以精确测量冥王星的质量，并确定其远小于最初的估计。\(^\text{[7]}\)其质量大约只有水星的二十分之一，使冥王星成为迄今最小的行星。尽管冥王星仍比小行星带中最大的天体谷神星大十倍以上，但其质量却只有地球月球的五分之一。\(^\text{[8]}\)此外，由于其具有一些异常特征，如较大的轨道离心率与较高的轨道倾角，愈发明显冥王星与其他行星类型不同。\(^\text{[9]}\)

在1990年代，天文学家开始在与冥王星相同的空间区域（即如今所称的柯伊伯带）以及更远处发现更多天体。\(^\text{[10]}\)其中许多与冥王星共享关键的轨道特征，冥王星开始被视为一类新天体的最大成员，即冥族（plutinos）。于是变得明显：要么这些较大的天体也需要被归类为行星，要么冥王星需要像谷神星在更多小行星被发现后那样被重新分类。\(^\text{[11]}\) 这使得一些天文学家不再将冥王星称为行星。多个术语包括“subplanet”（次行星）与“planetoid”（类行星体）等，被用于指代如今所谓的矮行星。\(^\text{[12][13]}\)天文学家也确信会有更多与冥王星大小相当的天体被发现，如果冥王星继续被视为行星，行星数量将会快速增长。\(^\text{[14]}\)

海王星外天体厄里斯（当时称为 2003 UB313）于 2005 年 1 月被发现；\(^\text{[15]}\)当时认为其尺寸略大于冥王星，因此一些报道非正式地称其为“第十颗行星”。\(^\text{[16]}\)因此，该问题在 2006 年 8 月的国际天文学联合会（IAU）大会上成为激烈争论的焦点。\(^\text{[17]}\)IAU 的最初草案提议将卡戎、厄里斯与谷神星列入行星名单。由于许多天文学家反对这一方案，乌拉圭天文学家 Julio Ángel Fernández 与 Gonzalo Tancredi 制定了另一项替代方案：他们提出设立一个中间类别，用于归类那些质量足以达到流体静力平衡（呈球形），但并未清除其轨道中微行星的天体。除了将卡戎从名单中移除之外，该新方案也将冥王星、谷神星与厄里斯排除在行星名单之外，因为它们均未清除其所在轨道。\(^\text{[18]}\)

尽管有人对此分类在系外行星中的适用性提出了担忧，\(^\text{[19]}\)该问题仍未得到解决；提议是等到能够观测到与矮行星体积相当的系外天体时再作决定。\(^\text{[18]}\)

在 IAU 给出“矮行星”定义的直接后续中，一些科学家表达了他们对 IAU 决议的不满。\(^\text{[20]}\)抗议活动包括汽车保险杠贴纸与 T 恤标语。\(^\text{[21]}\)厄里斯的发现者 Mike Brown 则赞同将行星数量减少至八颗的做法。\(^\text{[22]}\)

NASA 在 2006 年宣布将采用 IAU 制定的新指南。\(^\text{[23]}\)NASA 冥王星任务的负责人 Alan Stern 不认可当前 IAU 的行星定义，无论是将矮行星排除在行星类别之外，还是依据轨道特征（而非天体的固有特征）来定义矮行星。\(^\text{[24]}\)因此，他在 2011 年仍将冥王星称为行星，\(^\text{[25]}\)并将其他可能的矮行星，如谷神星与厄里斯，以及更大的卫星，也视作额外的行星。\(^\text{[26]}\)在 IAU 发布定义的数年前，他便使用轨道特征将“überplanets”（主导性的八颗行星）与“unterplanets”（矮行星）区分开来，并认为二者皆为“行星”。\(^\text{[27]}\)
\subsection{名称}
\begin{figure}[ht]
\centering
\includegraphics[width=6cm]{./figures/46cf07ad5d90f072.png}
\caption{欧拉图：展示国际天文学联合会（IAU）执行委员会对太阳系（不含太阳）内天体类别的构想} \label{fig_Dwarf_4}
\end{figure}
对大型亚行星天体的称呼包括 dwarf planet（矮行星）、planetoid（类行星，更泛化的术语）、meso-planet（中型行星，狭义上指介于水星与谷神星之间尺寸的天体）、quasi-planet（准行星），以及（在海王星外区域）plutoid（冥族天体）。然而，dwarf planet 这一术语最初是用于指代“最小的行星”，而不是“最大的亚行星”，并且至今仍以这种含义被许多行星地质学家使用。

dwarf planet 作为 asteroid（小行星）的同义词可以追溯至少到 1838 年。它也曾作为 giant planet（巨行星）的反义词，被用于包括类地行星——水星、金星、地球和火星——以及月球在内的天体。\(^\text{[28]}\)在 2006 年的分类中，IAU 规定矮行星不应被视为行星。对于 IAU 所界定的最大亚行星天体，其他不带有类似歧义或用法冲突的术语包括 quasi-planet\(^\text{[29]}\)与较早的 planetoid（“具有行星外形之物”）。\(^\text{[30]}\)Michael E.~Brown 指出，planetoid “是一个完全没问题的词汇，使用多年”，而用 dwarf planet 去指代“非行星”则是“很傻”（dumb），但这是因为 IAU 第三分部大会试图通过第二项决议让冥王星重新成为行星。\(^\text{[31]}\)事实上，5A 号决议草案原本将这些中等大小的天体称为 planetoids，\(^\text{[32][33]}\)但大会一致投票将其名称改为 dwarf planet。\(^\text{[34]}\)第二项决议 5B 则将矮行星定义为行星的一种子类型（与 Stern 最初的设想一致），与其它八颗行星区分为“classical planets”（经典行星）。在这种安排下，被否决的提案所列的 12 颗“行星”本可以通过“八颗经典行星 + 四颗矮行星”这一区分得以保留。然而，决议 5B 在通过 5A 的同一场会议中被否决了。\(^\text{[31]}\)由于决议 5B 未获通过，从语义上造成了“矮行星并非行星”这一不一致，因此出现了 nanoplanet、subplanet 等替代名词的讨论，但 IAU 的 CSBN 并未就更名达成共识。\(^\text{[35]}\)

在大多数语言中，“矮行星”都被更或少直译为对应词语：法语 planète naine、西班牙语 planeta enano、德语 Zwergplanet、俄语 karlikovaya planeta（карликовая планета）、阿拉伯语 kaukab qazm（كوكب قزم）、中文 矮行星、韩语 왜소행성／waesohangseong（矮小行星）或 왜행성／waehangseong（矮行星）。但在日语中，它们被称为 准惑星（junwakusei），意为“准行星”或“近似行星”（peneplanet，pene- 表示“几乎”）。

2006 年的 IAU 第 6a 号决议\(^\text{[36]}\)将冥王星认定为“新一类海王星外天体的原型”。该类别的名称与精确定性当时尚未指定，而是留待 IAU 日后确定；在该决议出台前的讨论中，这一类别的成员曾被称为 plutons 与 plutonian objects，但这两个名称均未被采纳，可能是因为地质学家反对这种命名会与地质学中的 pluton（深成岩体）混淆。\(^\text{[34]}\)

2008 年 6 月 11 日，IAU 执行委员会宣布了一个新术语 plutoid，并给出了定义：所有海王星外的矮行星皆为 plutoid。\(^\text{[37]}\)然而，IAU 的其他部门拒绝了这一术语：

……部分由于电子邮件沟通失误，WG-PSN（行星系统命名工作组）并未参与“plutoid”一词的选择。……事实上，该工作组在执行委员会会议之后所进行的投票否决了这一特定术语的使用……\(^\text{[35]}\)

“plutoid”这一类别继承了先前根据地质性质在“类地矮行星”谷神星与外太阳系“冰矮星”之间所做的区分。\(^\text{[38]}\)这一区分是三分太阳系概念的一部分：即太阳系由内侧类地行星、中部巨行星、以及外侧冰矮星构成，而冥王星是该外侧组的代表成员。\(^\text{[39]}\)“ice dwarf（冰矮星）”过去也曾用作海王星外所有小行星（TNOs）的总称，或指外太阳系的“冰质小行星”；其中一种尝试性的定义是：冰矮星“比典型彗核更大，且比典型小行星更富含冰”。\(^\text{[40]}\)

自 Dawn 号任务以来，人们已经认识到谷神星是一个地质上富冰的天体，并可能起源于外太阳系。\(^\text{[41][42]}\)因此谷神星也被称作冰矮星。\(^\text{[43]}\)\subsection{分类标准}
\begin{figure}[ht]
\centering
\includegraphics[width=10cm]{./figures/cf2087e9444efd75.png}
\caption{} \label{fig_Dwarf_5}
\end{figure}
“矮行星”这一类别的出现，源于关于“何谓行星”这一概念究竟应依据动力学标准还是地球物理学标准而产生的冲突。就太阳系动力学而言，最重要的区分在于：一类天体能够在其轨道邻域内表现出引力支配地位（从水星到海王星），而另一类天体则不能（例如小行星与柯伊伯带天体）。一个天体在其质量达到某一程度时，其地幔会在自身重量作用下变得具可塑性，从而使天体获得圆形外貌，此时它可能具备动力学意义上的“行星形”地质活动。由于获得圆形所需的质量远低于实现轨道邻域引力清空所需的质量，因此存在这样一类天体：其质量足以具备“类世界的外貌”以及类似行星的地质活动，但不足以清除其轨道邻域。例如小行星带中的谷神星，以及柯伊伯带中的冥王星。[47]

动力学研究者通常倾向将“引力支配轨道邻域”作为行星的判定标准，因为从他们的角度看，将质量较小的天体与其邻域群体归为一类更为合理，例如将谷神星视为一颗较大的小行星，将冥王星视为一颗较大的柯伊伯带天体。[48][49] 地球科学家则通常以“是否达到流体静力平衡形状”作为判定阈值，因为从他们的视角来看，像谷神星这样拥有内部驱动地质活动的天体，与拥有行星地质的经典行星（如火星）更为相似，而与缺乏内部地质驱动的小型小行星并不相同。正因为存在这一中间层级，才有必要设立“矮行星”这一类别以描述该类型天体。[47]
\subsubsection{轨道支配性}
Alan Stern 与 Harold F. Levison 在 2000 年提出了参数 Λ（大写 lambda），用以表示一次遭遇导致轨道发生给定偏转的概率。[27] 在 Stern 的模型中，该参数的取值与天体质量的平方成正比，并与轨道周期成反比。该值可用于估计天体清除其轨道邻域的能力，其中  Λ > 1 表示该天体最终能够清除其轨道区域。研究发现，从最小的类地行星到最大的一批小行星与柯伊伯带天体之间，Λ 的数值存在五个数量级的差距。[44]

利用这一参数，Steven Soter 以及其他天文学家主张依据天体“无法清除其轨道邻域”来区分行星与矮行星：行星能够通过碰撞、俘获或引力扰动清除其轨道附近的小型天体（或建立轨道共振以防止碰撞）；而矮行星由于质量不足，无法完成这一过程。[27]Soter 进一步提出了一个称为“行星判别式”的参数，记作 µ（mu），其代表轨道区域实际清洁度的实验度量（µ 的计算方式为：候选天体的质量除以与其共享轨道区的所有其他天体的总质量），其中  µ > 100 被视为已清除轨道邻域。[44]

Jean-Luc Margot 对 Stern 与 Levison 的概念进行了改进，提出了一个类似的参数 Π（大写 Pi）。[46]该参数基于理论推导，避免了 Λ 中所依赖的经验性数据。Π > 1 即被视为行星，同样地，行星与矮行星之间在 Π 的取值上也存在若干个数量级的差距。

此外，还有其他一些方案试图区分行星与矮行星，[20] 但 2006 年 IAU 的定义采用的正是这一概念。[34]
\subsubsection{流体静力平衡}
当天体的自引力产生足够的内部压力时，其物质会进入塑性状态，而足够的塑性会使高起伏下沉、低洼区域被填平，这一过程称为重力松弛（gravitational relaxation）。**
直径小于数公里的天体主要受非重力因素主导，往往呈不规则形状，并可能是碎石堆（rubble piles）。更大的天体中，引力开始变得重要但尚未完全占优势，此时它们呈现出“马铃薯状”。随着天体质量的增加，其内部压力进一步升高，结构更为坚固，形状也随之更加圆滑，直到内部压力足以压过其抗压强度，使天体进入流体静力平衡（hydrostatic equilibrium）。此时，天体在其自转和潮汐效应允许的条件下，已圆整到极限，呈现椭球体形状。这就是“矮行星”的判定界限。[50]

如果一个天体处于流体静力平衡，那么其表层若覆盖一层全球性的液体，该液体的表面形状（除去陨坑裂隙等小尺度地形）将与该天体整体的形状一致。不自转的天体将呈球形，而自转天体呈椭球形。自转越快，天体越扁甚至可能呈不等轴形。若此类自转天体被加热直至熔融，其形状也不会改变。一个极端的例子是 Haumea，它因快速自转可能呈不等轴形，其最长轴长度约为极半径的两倍。

如果天体拥有一个质量巨大的近伴，则潮汐作用会逐渐减慢其自转，直至其达到潮汐锁定（tidally locked）；即始终以同一面朝向其伴星。潮汐锁定的天体也呈不等轴形，但有时仅有轻微差异。地球的月球处于潮汐锁定状态，所有气体巨行星的圆形卫星亦如此。冥王星与其卫星卡戎彼此潮汐锁定；阋神星（Eris）与其卫星迪斯诺米亚（Dysnomia）亦如此；Orcus 与 Vanth 很可能也处于相同状态。[citation needed]

矮行星并无特定的大小或质量上限，因为这些并非其定义依据。实际上也不存在明确的上限：例如，在太阳系的极外区域，一个质量超过水星但尚未清除其轨道区的天体，也将被归类为矮行星而非行星。事实上，Mike Brown 曾试图寻找这样的天体。[51]

下限由天体达到并保持流体静力平衡的要求决定，但具体大小或质量阈值取决于其成分与热演化历史，而非简单的质量。IAU 2006 年新闻稿[52] 的问答部分估计，质量超过 (0.5\times 10^{21},\mathrm{kg}) 且半径大于 (400,\mathrm{km}) 的天体“通常”会处于流体静力平衡状态（形状……通常由自引力决定），但所有边缘个案都必须通过观测来确定。[52]

这一估计与截至 2019 年对海王星外完整固体天体的判断相近，其中 Salacia（(r = 423\pm 11,\mathrm{km})，(m = (0.492\pm 0.007)\times 10^{21},\mathrm{kg})）在 2006 年问答标准与后续研究中均被视为边界案例；而 Orcus 稍高于该预期上限。[53] 就已测质量的天体而言，没有其他对象接近上述质量上限，但一些未测质量的天体接近所预期的尺寸下限。[citation needed]
